\section{Conditional screening inputs}\label{sec:inputs}

We define one family of screened conditional fields and state the
foundation estimates used to study it.  Their complete proofs are in
\cite{OuterRadius}; Appendix~\ref{sec:interface} collects the parameter
order and provider map.  All constants and thresholds below are uniform
over neutral ground states.

\subsection{Energy offset and count bounds}

For a fixed nuclear charge $Z$, let $E_n=E_n(Z)$ be the bottom of the
$n$-electron Hamiltonian with the same kinetic and spin conventions as
\eqref{intro:hamiltonian}, and put $E=\inf_{n\ge0}E_n$.
A normalized state has offset $\delta$ if its energy is at most $E+\delta$.

\begin{proposition}[Uniform offset and raw counts]\label{inp:energy-counts}
There are constants $D_0,Z_0<\infty$ such that every normalized neutral
ground state with $Z\ge Z_0$ has offset at most $D_0$.
For every fixed $0<\alpha<\beta<\infty$ and every $u>0$, its raw
electron configuration satisfies
\begin{equation}\label{inp:counts}
 \begin{aligned}
 \left\|\#\{i:\alpha u<|x_i|<\beta u\}\right\|_{L^2(\Prob)}
 &\le C_{\alpha,\beta}
       \bigl(\max\{u^{-3},1\}+\sqrt{D_0u}\bigr),\\
 \E\#\{i:|x_i|>1\}&\le C(1+\sqrt{D_0}).
 \end{aligned}
\end{equation}
\end{proposition}

\begin{proof}
Monotonicity of the sector energies
\cite[\FoundationMonotonicity]{OuterRadius} and the bounded offset of the
$(Z-1)$-electron sector \cite[\FoundationOffset]{OuterRadius} give
$E\le E_Z\le E_{Z-1}\le E+C$.  Take $D_0=C$.
The first assertion in \eqref{inp:counts} is
\cite[\FoundationCounts]{OuterRadius}.  The second follows from
\cite[\FoundationDeficit]{OuterRadius} with $t=4$, since neutrality gives
\[
 \E\#\{i:|x_i|>1\}=Z-\E\#\{i:|x_i|<1\}.
\]
Boundary spheres have zero probability, so either convention for
annular endpoints gives the same estimates.
\end{proof}

\subsection{One family of observations and packets}

The comparison estimate concerns a smoothed density conditioned on noisy
position observations.  As the observation index increases, we discard
finer observations; the propagation construction uses conditional
averaging over the discarded data.  We use one packet family throughout,
so that comparison and propagation concern the same fields.

Fix a real smooth radial function $g$, supported in the unit ball and
normalized by $\int g^2=1$, and set
$g_t(y)=t^{-3/2}g(y/t)$.  Fix the universal constants
\[
 w=10^{-5},\qquad L=10^5,\qquad 0<c_1<(10L)^{-1}.
\]
Also fix $\eps>0$ sufficiently small for the barrier construction below.
For $0<s<1$ and $Z$ large enough that $2\eps Z^{-1/3}\le s$, define
\begin{equation}\label{inp:observations}
 r_0=\eps Z^{-1/3},\qquad r_j=2^jr_0,\qquad
 J=\max\{j:r_j\le s\},\qquad s/2<r_J\le s.
\end{equation}
For each $j<J$ observe the unordered array
$\{x_i+r_j^{1.01}U_{ji}:1\le i\le Z\}$.  Every Cartesian
coordinate of every $U_{ji}$ is independent, independent of the raw
configuration, and has density
$c_\zeta\exp[-1/(1-u^2)]\mathbf1_{\{|u|<1\}}$.
Let $\mathcal F_j$ retain the arrays with indices $j,\ldots,J-1$;
let $\mathcal F_J$ be trivial.  Thus the sigma-fields decrease with $j$.

For a possible electron center $x$, define the master packets and their
conditional densities by
\begin{equation}\label{inp:packets}
 \begin{gathered}
 d_x=|x|,\qquad \widehat d_x=\max\{d_x,r_0\},\qquad
 t_x=c_1\widehat d_x\min\{\widehat d_x,s\}^{w},\\
 \mathcal K_x(y)=g_{t_x}(y-x)^2,\qquad
 \mu_j(y)=\E\left[\sum_{i=1}^Z\mathcal K_{x_i}(y)
                  \,\middle|\,\mathcal F_j\right],
 \qquad H_j=V-K*\mu_j .
 \end{gathered}
\end{equation}
The upper clamp acts only on the factor raised to $w$.
In particular $t_x\ge c_1r_0^{1+w}>0$, every packet has mass one,
and its support is contained in $B(x,t_x)$.
Unlike the unconditional raw density $\rho_{\Psi_Z}$, the density
$\mu_j$ depends on the observed data and includes packet smoothing.
Section~\ref{sec:tails} removes both operations to recover the expected
raw mass defining the radius.

We use the conditional-integral versions in
\cite[\FoundationObservations]{OuterRadius}: at data of positive marginal density,
integrate against the raw law multiplied by the observation likelihood
and divided by that marginal density.  These versions are jointly
measurable.  For each finite system, every spatial derivative of
$\mu_j$ has a deterministic bound, since the widths have a positive
lower bound.  Almost surely,
\begin{equation}\label{inp:mass-poisson}
 \mu_j\ge0,\qquad \int\mu_j=Z,\qquad
 \Delta H_j=-4\pi Z\delta_0+4\pi\mu_j .
\end{equation}
The remainder $H_j-V$ is continuous and locally Lipschitz even at zero.
For example, splitting the Coulomb potential and its gradient at unit
distance uses boundedness of $\mu_j$ near the singularity and its finite
mass farther away.  The corresponding regularity bounds may depend on
$Z$.  The conditional tower identity for $\mu_j$ holds pointwise in
space; the identity for $H_j$ holds away from the nucleus, and that
for the continuous extension of $H_j-V$ holds everywhere.  Both
identities hold after integration against smooth compactly supported
tests.

\begin{samepage}
\subsection{The inverse comparison}

\begin{proposition}[Conditional comparison]
\label{inp:comparison-prop}
There is a universal $C_{\mathrm{cap}}$ with the following property.
Fix $L_1\ge2$, $0<h_l<h_h<\infty$, and $0<\xi<h_l/16$.
For all sufficiently small $s>0$, depending on these parameters and
$D_0$, and each $j<J$, there is an event in $\mathcal F_j$ whose
complement has probability at most $Cr_j^{25}$, on which, simultaneously
for $r_j\le |y|\le4L_1r_j$ and $h_l\le h\le h_h$,
\begin{equation}\label{inp:comparison}
 \begin{gathered}
 H_j(y)\le C_{\mathrm{cap}}|y|^{-4},\\
 |y|^6\mu_j(y)>kh^{3/2}
       \quad\Longrightarrow\quad |y|^4H_j(y)\ge h-\xi,\\
 |y|^6\mu_j(y)<kh^{3/2}
       \quad\Longrightarrow\quad |y|^4H_j(y)\le h+\xi.
 \end{gathered}
\end{equation}
The probability constant may depend on the fixed comparison
parameters, but $C_{\mathrm{cap}}$ does not.
The estimates are uniform in $Z$, the ground state, $j$, and $J$.
\end{proposition}
\end{samepage}

This is \cite[\FoundationComparison]{OuterRadius}.
The observations and packets in \eqref{inp:observations}--\eqref{inp:packets}
are chosen before $L_1,h_l,h_h,\xi$.  Varying the comparison accuracy or
its annular width therefore does not change the conditional fields.
This common choice is essential when we combine increasingly accurate
comparisons with a single barrier construction.

\subsection{Intermediate subsolutions}

\begin{proposition}[Intermediate subsolutions]
\label{inp:barrier-prop}
For the fixed constants above, there are $B>0$, $C_{\mathrm{inv}}<\infty$,
$s_{\mathrm{bar}}>0$, and a finite integer threshold
$Z_{\mathrm{bar}}(s)$ for every $0<s\le s_{\mathrm{bar}}$, such that
the following holds.  For every $Z\ge Z_{\mathrm{bar}}(s)$ and
every neutral ground state,
\cite[\FoundationIntermediate]{OuterRadius} supplies, at each $0\le j\le J$,
$\mathcal F_j$-measurable functions $u_j,p_j$ satisfying
\begin{equation}\label{inp:barriers}
 \begin{gathered}
 \Delta u_j\ge-4\pi Z\delta_0+
     4\pi\mathbf1_{\{|y|<r_j\}}\mu_j-4\pi p_j+
     \mathbf1_{\{|y|\ge r_j\}}f(u_j)
       \quad\hbox{in }\mathcal D'(\R^3),\\
 B|y|^{-4}\left(1-\frac{r_j}{8|y|}\right)
       \le u_j(y)\le C_{\mathrm{inv}}|y|^{-4}
       \quad (|y|\ge r_j),\\
 p_j\ge0,\qquad \supp p_j\subset\{|y|\le r_j\},\qquad
 \E\int p_j\le C(s^{32}+s^\gamma),\qquad
 \gamma=\frac{19}{2}-3w>0.
 \end{gathered}
\end{equation}
For each finite system and index, $p_j$ is bounded and $u_j-V$ is
continuous and locally Lipschitz through zero.  The constants $B$,
$C_{\mathrm{inv}}$, and $C$ and the thresholds are independent of the
ground state and of the number of dyadic scales.
We enlarge $Z_{\mathrm{bar}}(s)$ to include $Z\ge Z_0$ and $2r_0\le s$.
\end{proposition}

The bound on the expected error is valid at every intermediate index.
Its persistence is important enough to recall: the propagation first
replaces $p_j$ by
\[
 \widetilde p=p_j+
   \mathbf1_{\{\text{bad data at step }j\}}
   \mathbf1_{\{r_j\le|y|<2r_j\}}\mu_j
\]
and then sets
$p_{j+1}=\E[\widetilde p\mid\mathcal F_{j+1}]$;
as in the proof of \cite[\FoundationIntermediate]{OuterRadius}.
The added source is nonnegative and conditional averaging preserves its
expected integral.  Hence $\E\int p_j$ is nondecreasing in $j$,
so the final bound controls all preceding indices.
We fix all parameters of this construction once.  The only parameters
varied later are those of \cref{inp:comparison-prop}, applied to the same
observations and packets.

At a fixed index $j$, the identities, finite-system regularity, and
almost-sure subsolution properties may be placed on one
$\mathcal F_j$-measurable set of full measure.  Continuity and a countable
dense family of compactly supported distributional tests justify the
simultaneous spatial assertions.  We fix these versions before selecting
any observation data.  The comparison inequalities of
\cref{inp:comparison-prop} hold on their separate high-probability event,
which will be intersected with this full-measure set.
